% Page 006

\seite{6}

Sonnt sich an der Tempelwand\ob
Wo geflammt der Opferbrand."

\margmark{Das hohe Alter\\Sefferweichs}%
Aber Sefferweich ist alt, sehr alt, das beweisen\ob
schon seine ältesten Namen: Wrivic, Uwriche, Orwic,\ob
Orwich, Orweiche. Seine Entstehung verliert sich ins\ob
Dunkel der Zeiten. Daß die Kelten sich hier nieder\ohb
gelassen haben, ist über allen Zweifel erhaben;\ob
\margmark{etwa um 1840}%
denn vor mehreren Jahren fand man in der Gegend\ob
von Neidenbach einen viereckigen bearbeiteten\ob
Stein beim Pflügen, mit einer Inschrift, welche\ob
von Altertumsforschern als eine keltische erkannt\ob
worden sein soll. Dieser Stein wurde\ob
\margmark{Funde aus der\\Zeit der Kelten}%
dem Finder Philipp Roths aus Neidenbach\ob
mit 25 Thaler bezahlt und befindet sich gegenwärtig\ob
in Bonn. Auch fand man im Jahre 1852 auf\ob
dem selben Wege von Sefferweich nach Staffelstein\ob
auf der Höhe rechts vom Wege in einem Acker\ohb
felde ein keltisches Grab. Leider ist dasselbe aus\ob
Unkenntnis des Finders, der einen großen Schatz\ob
zu finden meinte, zerstört worden. Der Inhalt\ob
desselben, bestehend aus ungewöhnlich großen Knochen\ob
und aus anderm Unnützen ist weggeworfen wor\ohb
den. Die Möglichkeit bleibt nicht ausgeschlossen, daß\ob
es einem Späteren gelingen kann, noch andere\ob
Gräber auf derselben Flur zu finden. Haben\ob
wir auch, wie leicht erklärlich, keine bündigen
\pend
